X3DH is built upon Elliptic Curve Diffie-Hellman (ECDH) key exchange. The security of ECC is derived from the Elliptic Curve Discrete Logarithm Problem (ECDLP), which states that given a point multiplication:

$$
Q = kP
$$

it is computationally infeasible to determine the scalar \( k \) given only \( Q \) and the base point \( P \).

An elliptic curve over a finite field \( \mathbb{GF}(p) \) is defined as:

$$
y^2 = x^3 + ax + b
$$

where \( a \) and \( b \) define the curve parameters, and \( p \) is a prime number, while a base point, or generator, P is a point on the elliptic curve which forms a cyclic subgroup. In a cyclic subgroup, every other element (point) in the subgroup can be written as the point multiplication between a scalar, k, and P, which represents                                                                                                                                             adding the point P to itself k times. The order of P, denoted \( \text{ord}(P) \), is the smallest integer n such that nP=\( \infty \), where \( \infty \) acts as the identity in the elliptic curve group (\( \infty \)+P=P).
In ECDH key exchanges, Alice and Bob both choose a random private key, \( a \) and \( b \) (key clamping ensures that this private key isn't too large or small), from the order and calculate a public key \( A = aP \) and \( B = bP \), which they exchange with each other. Afterwards, they use their individual private keys, which any adversary would have no knowledge of, to calculate a shared key \( K = bA = aB = abP \). Only relying on one key exchange has some drawbacks, particularly being vulnerable to man in the middle attacks. Additionally in order to achieve forward and backward secrecy, a new ECDH exchange needs to be performed each time, requiring both users to be online to generate a new public and shared key. Otherwise, any adversary who manages to compromise a single message would also gain access to all previous/future messages.